% Fichier généré automatiquement par tools/generate_corrections.py.
% Source : PERF/PERF-06-Precision/63_BancHydraulique/63_BancHydraulique.tex
% Statut : complet (5/5 question(s) renseignée(s), 0 reprise(s)).
\begin{corrigebox}[Corrigé extrait de la banque]
\CorrigeQuestion{1}
Le débit de fuite est nul; donc $\Delta Q_e(p)=0$.

\textbf{Cas 1 : cours sur la précision connu \textit{-- Attention à avoir le même type d'entrée/sortie}}

La FTBO est de classe nulle ($C(p)$ est un gain, $H_{\text{pom}} (p)$ et $H_{\text{pre}} (p)$ de classe 0). Le gain de la Boucle ouverte est $K_{\text{BO}}=K_p K_m K_{\text{pom}}K_{\text{cap}}$.

Si l'entrée est un échelon d'amplitude $P_0$, l'écart statique est donc donné par 
$\varepsilon_S = \dfrac{P_0}{1+K_{\text{BO}}}= \dfrac{P_0}{1+K_p K_m K_{\text{pom}}K_{\text{cap}}}$.

\textbf{Cas 2 : cours sur la précision peu connu -- À savoir faire, mais on perd un peu de temps... \textit{-- Attention à avoir le même type d'entrée/sortie}}
Si on connait quand même un petit peu son cours, on a 
$\varepsilon(p)
=
\dfrac{P_{\text{con}}(p)}
{1+K_P \dfrac{K_{\text{pom}}}{1+T_2 p} \dfrac{K_m}{1+T_1 p} K_{\text{cap}}}$.

On a alors,
$\varepsilon_s = \lim\limits_{p\to 0} p 
\dfrac{\dfrac{P_0}{p}}
{1+K_P \dfrac{K_{\text{pom}}}{1+T_2 p} \dfrac{K_m}{1+T_1 p} K_{\text{cap}}}$
$=
\dfrac{P_0}
{1+K_P K_{\text{pom}}K_m K_{\text{cap}}}
$

\textbf{Cas 3 : cours sur la précision pas connu -- À savoir faire, mais on perd beaucoup peu de temps...}

En utilisant la formule de Black, on a $P_e(p) 
= P_{\text{con}}(p) K_{\text{cap}} 
\dfrac{K_P \dfrac{K_{\text{pom}}}{1+T_2 p} \dfrac{K_m}{1+T_1 p} }
{1+K_P \dfrac{K_{\text{pom}}}{1+T_2 p} \dfrac{K_m}{1+T_1 p} K_{\text{cap}}}$

$= P_{\text{con}}(p) K_{\text{cap}}(p) 
\dfrac{K_P K_{\text{pom}}K_m }
{\left(1+T_2 p\right)\left(1+T_1 p\right)+K_P K_{\text{pom}} K_m K_{\text{cap}}}$

En passant à la valeur finale avec une entrée échelon, on a 
$\lim\limits_{t\to +\infty}P_e(t)$ $ =  P_{0} K_{\text{cap}} 
\dfrac{K_P K_{\text{pom}}K_m }
{1+K_P K_{\text{pom}} K_m K_{\text{cap}}}$

L'écart statique est donc donné par
$
\varepsilon_S = P_0 - P_{0} 
\dfrac{K_P K_{\text{pom}}K_m K_{\text{cap}} }
{1+K_P K_{\text{pom}} K_m K_{\text{cap}}}$
$= P_0
\dfrac{1+K_P K_{\text{pom}} K_m K_{\text{cap}} - K_P K_{\text{pom}}K_m K_{\text{cap}} }{1+K_P K_{\text{pom}} K_m K_{\text{cap}}}
$

$= 
\dfrac{P_0}{1+K_P K_{\text{pom}} K_m K_{\text{cap}}}$
\CorrigeQuestion{2}
On souhaite que l'écart statique soit inférieure à 5\% soit 0,05 pour une entrée unitaire. 

On cherche donc $K_P$ tel que 
$\dfrac{1}{1+K_P K_{\text{pom}} K_m K_{\text{cap}}}<0,05 $
$ \Leftrightarrow 1<0,05\left(1+K_P K_{\text{pom}} K_m K_{\text{cap}}\right)$

$ \Leftrightarrow \dfrac{1 - 0,05}{0,05 K_{\text{pom}} K_m K_{\text{cap}}}<K_P $

Soit $ K_P > \dfrac{1 - 0,05}{0,05 \times 1,234 \times 10^7\times 3,24 \times  2,5 \times 10^{-8}}  \Rightarrow K_P > 19$.
\CorrigeQuestion{3}
Dans ce cas il n'y a pas d'intégrateur avant la perturbation échelon. Il faut savoir faire le calcul.

On peut utiliser la << lecture directe >> :
$P_e(p)= P_r(p)\indice{H}{pre} - \Delta Q_e(p) \indice{H}{fui}(p)$

$= \indice{H}{pre}(p)\indice{H}{pom}(p) C(p) \varepsilon(p)- \Delta Q_e(p) \indice{H}{fui}(p)$

$= -\indice{H}{pre}(p)\indice{H}{pom}(p) C(p) \indice{K}{cap} P_e(p)- \Delta Q_e(p) \indice{H}{fui}(p)$.

$\Leftrightarrow P_e(p) \left(1+\indice{H}{pre}(p)\indice{H}{pom}(p) C(p) \indice{K}{cap} \right)
=- \Delta Q_e(p) \indice{H}{fui}(p)$

$\Leftrightarrow \dfrac{P_e(p)}{\Delta Q_e(p)} 
=-  \dfrac{\indice{H}{fui}(p)}{1+\indice{H}{pre}(p)\indice{H}{pom}(p) C(p) \indice{K}{cap}}$

Calculons $\indice{\varepsilon}{pert}(p) = -  \dfrac{\indice{H}{fui}(p)}{1+\indice{H}{pre}(p)\indice{H}{pom}(p) C(p) \indice{K}{cap}} \Delta Q_e(p) \indice{K}{cap}$.

On a alors $\indice{\varepsilon}{pert}= \tvf{\varepsilon}$ 
$= \lim\limits_{p\to 0} - p\times  \dfrac{\indice{H}{fui}(p)}{1+\indice{H}{pre}(p)\indice{H}{pom}(p) C(p) \indice{K}{cap}} \dfrac{\Delta Q_0}{p} \indice{K}{cap}$

$= - \dfrac{K_f\Delta Q_0 \indice{K}{cap}}{1+K_m\indice{K}{pom} K_P \indice{K}{cap}} $
\CorrigeQuestion{4}
Pour $\Delta Q_e = \SI{5e-4}{m^3.s^{-1}}$, il faut $\indice{\varepsilon}{pert}<40\times 10^5$ (Pa) soit

$\dfrac{K_f\Delta Q_0 \indice{K}{cap} }{1+K_m\indice{K}{pom} K_P \indice{K}{cap}}  < 40\times 10^5$
$\Rightarrow  K_f\Delta Q_0 \indice{K}{cap}  < 40\times 10^5\left(1+K_m\indice{K}{pom} K_P \indice{K}{cap}\right)$
$\Rightarrow \dfrac{K_f\Delta Q_0 \indice{K}{cap}  -40\times 10^5}{40\times 10^5 K_m\indice{K}{pom} \indice{K}{cap} }< K_P $
$\Rightarrow K_P > -1$
\CorrigeQuestion{5}
Je vous laisse faire le calcul... Il faut savoir le faire le plus vite possible.
Il faut d'abord calculer la FTBF, la mettre sous forme canonique, déterminer 
$\indice{\xi}{BF}=\dfrac{T_1+T_2}{2\sqrt{T_1T\left( 1+K_P K_M \indice{K}{Pom}\indice{K}{Cap}\right)}}$ puis
determiner $K_P$ tel que $\indice{\xi}{BF}=1$.

\begin{enumerate}
  \item $\varepsilon_{\text{con \%}} = \dfrac{1}{1+K_PK_m K_{\text{pom}} K_{\text{cap}} }$;
  \item $K_P > 19$;
  \item $\varepsilon_{\text{pert}} = \Delta Q_e \dfrac{K_f}{1+K_{\text{cap}}K_PK_mK_{\text{pom}}}$;
  \item $K_P > -1$.
  \item $K_P < 0,125$. Il est impossible de vérifier les trois conditions avec un correcteur proportionnel.
\end{enumerate}
\end{corrigebox}
